% Created 2021-04-22 Thu 17:31
% Intended LaTeX compiler: pdflatex
\documentclass[9pt, b5paper]{article}
\usepackage[UTF8]{ctex}
\usepackage{fontspec}
\usepackage{graphicx}
\usepackage{xcolor}
\usepackage{multirow}
\usepackage{multicol}
\usepackage{float}
\usepackage{textcomp}
\usepackage{geometry}
\geometry{left=1.2cm,right=1.2cm,top=1.5cm,bottom=1.2cm}
\usepackage{algorithm}
\usepackage{algorithmic}
\usepackage{latexsym}
\usepackage{natbib}
\usepackage{listings}
\usepackage{minted}
\usepackage[xetex,colorlinks=true,CJKbookmarks=true,linkcolor=blue,urlcolor=blue,menucolor=blue]{hyperref}
\author{deepwaterooo}
\date{\today}
\title{成长的故事 —— 我和表哥}
\hypersetup{
 pdfauthor={deepwaterooo},
 pdftitle={成长的故事 —— 我和表哥},
 pdfkeywords={},
 pdfsubject={},
 pdfcreator={Emacs 27.1 (Org mode 9.3)}, 
 pdflang={English}}
\begin{document}

\maketitle
\tableofcontents


\section{小弱弱躲猫猫记（2）: 躲猫猫般的中学生涯}
\label{sec:org9383949}

\begin{center}
\includegraphics[width=.9\linewidth]{./pic/backups_plans_20210422_101254.png}
\end{center}

等到大姐到了相亲、谈恋爱的年龄，我就感觉爸妈更喜欢大姐！

\begin{center}
\includegraphics[width=.9\linewidth]{./pic/backups_plans_20210422_095031.png}
\end{center}

这个孩子的成长，几乎是一种完全放养型的成长。她的成长过程中，很多观念没有树立起来，她很孤单、没什么力量。 

\begin{center}
\includegraphics[width=.9\linewidth]{./pic/backups_plans_20210422_094911.png}
\end{center}

当一个孩子主动与父母谈心，聊及某事的时候，当妈妈的当初也没有能够注意到我的反常，没有追问我是否有什么事情瞒着她。感觉那个时候，自己与父母之间早已经有着深深的隔痕。

\begin{center}
\includegraphics[width=.9\linewidth]{./pic/backups_plans_20210422_101554.png}
\end{center}

这才是作为晚辈、更应该是作为父母最大的悲哀吧：那个时候，我都没有觉得父母应该是自己最值得信任的人、没有这样一个概念！！！完全没有那份应有的信任。

\begin{center}
\includegraphics[width=.9\linewidth]{./pic/backups_plans_20210422_094753.png}
\end{center}

当初一我有一个还不错、很温暖的女生朋友的时候，我试探性地与她探讨过自己心里的难题。

\begin{center}
\includegraphics[width=.9\linewidth]{./pic/backups_plans_20210422_101717.png}
\end{center}

小学毕业事件后，我认为自己心里开始沉淀更多的自卑，如果小学时候不能加入、不能参与到自已喜欢的青梅竹马队列共同讨论电视剧、不常说话也算作是一种最初的自卑的话。

那个时候的自己：自认为自己自卑吗，还是没有意识呢？我们再寻找一下、找一找当年的那个自己和当时的状态，至少是在2011年11月左右回忆起来的时候自己所认定的状态。

\begin{center}
\includegraphics[width=.9\linewidth]{./pic/backups_plans_20210422_101921.png}
\end{center}

\begin{center}
\includegraphics[width=.9\linewidth]{./pic/backups_plans_20210422_102045.png}
\end{center}

\begin{center}
\includegraphics[width=.9\linewidth]{./pic/backups_plans_20210422_102136.png}
\end{center}

那时的自己，学习上确实是有信心的，后来的中考到县城里去考，第一次大型考试，也不曾紧张。 

\begin{center}
\includegraphics[width=.9\linewidth]{./pic/backups_plans_20210422_102214.png}
\end{center}

\begin{center}
\includegraphics[width=.9\linewidth]{./pic/backups_plans_20210422_102229.png}
\end{center}

在当时初三学生的学习竞赛里，全年级只有我一个人两门课得到复试的机会，最终也只有我一个人拿到《化学》单科省三等奖的成绩，其它两个与我一起复试过的小伙伴成绩都没有出线。 

\begin{center}
\includegraphics[width=.9\linewidth]{./pic/backups_plans_20210422_102326.png}
\end{center}

从小喜欢数学、数学单科的成绩也非常好，几乎没错什么。

\begin{center}
\includegraphics[width=.9\linewidth]{./pic/backups_plans_20210422_101951.png}
\end{center}

这应该是初中生、性格还没有形成、不够定性的时候的心思吧，是一种执念：认为自己学习好，就可以渺视体育！

所谓恃才傲物之类的性格也是如此形成的吗？

谁应该是那个为这种少儿执念、错误观念负责任的人？家长、老师（体育老师？班主任老师？），我自己？

\begin{center}
\includegraphics[width=.9\linewidth]{./pic/backups_plans_20210422_102437.png}
\end{center}

\begin{center}
\includegraphics[width=.9\linewidth]{./pic/backups_plans_20210422_102457.png}
\end{center}

那时的自己有那么点儿，不懂生活中的事的调调！

\begin{center}
\includegraphics[width=.9\linewidth]{./pic/backups_plans_20210422_102534.png}
\end{center}

比如一个细节吧。我进到办公室，明明知道化学老师谈论的就是我呀，我是我就是不生气呀，根本没把他的谈话当回事儿！

\begin{center}
\includegraphics[width=.9\linewidth]{./pic/backups_plans_20210422_171112.png}
\end{center}

这个平民家的女儿，其实还是比较感恩的！但生活中的事，也确实有点儿傻傻的了？！

\begin{center}
\includegraphics[width=.9\linewidth]{./pic/backups_plans_20210422_102629.png}
\end{center}

初三数学老师兼班主任、曾经以她自己的爱情、间接为我树立过的爱情观。

\begin{center}
\includegraphics[width=.9\linewidth]{./pic/backups_plans_20210422_171522.png}
\end{center}

高一高二的我：不是说刚过去的初三学习已经很好，已经比较自信了吗？为什么过了一个暑假就又变回原型了呢？

我个人的理解是觉得，在我“浸泡、沉浸式”小学初中艰难漫长的成长环境里，仅只一年的高三短暂自信，有一种人格不稳定、自信心不稳定。自信心积累得时间太短，显得不足够强大，过目即忘、转瞬即逝！

\begin{center}
\includegraphics[width=.9\linewidth]{./pic/readme_20210422_172512.png}
\end{center}

看了这所有的过往、与精神力量的分析，再来看我认舅舅时，“本能地‘冲’了上去”这种本能和舅舅那时带给过我的力量，就不难理解我会对一个遥远的国度有向往了吧！

后来高二高三的结果，之前的陈述里已经有很多了，这里关于高二在国内认了舅舅、高三高考意外等这里暂且略过。

\section{小弱弱躲猫猫记（3）: 大学班长（属羊水瓶座）}
\label{sec:orge05e709}
\section{小弱弱躲猫猫记（4）: 假如生活在那时转折}
\label{sec:orgec017f1}

\section{小弱弱躲猫猫记（1）: 青梅竹马（属马白羊座）}
\label{sec:orgfa8606e}

我上小学前，在当时家庭环境下，放过两年的牛。那时候，伯伯家的堂哥（比我大姐大一岁，比我大9岁）常常会帮忙照看着那时幼小的我。

\begin{center}
\includegraphics[width=.9\linewidth]{./pic/backups_plans_20210422_095907.png}
\end{center}

我的青梅竹马，是邻居村庄上一个同姓的侄儿。比我大一岁，我上小学一年级时，他与他们村另外两个男生一起留级到我们班，并成为班长第一批三个少先队员之一。没有猜错的话，他应该是属马白羊座的男孩。

\begin{center}
\includegraphics[width=.9\linewidth]{./pic/backups_plans_20210422_074655.png}
\end{center}

他很有亲和力（一如后来俺们的亲爱的表哥，我表哥给人的感觉更为亲切亲密，我们像是在哪里见过，眸子里仿佛能看见看透对方的三生三世！），也很有领导精神，能把整个村、邻村的同学、大小伙伴全都团结在一起！

\begin{center}
\includegraphics[width=.9\linewidth]{./pic/backups_plans_20210422_094617.png}
\end{center}

小时候的自己有一段时间，是很好奇自己从怎么、从哪里来到这个世界上的，所以也经常拿这个问题问妈妈。妈妈总说，我是爸爸上山砍柴时，从山上捡回来的娃娃！

\begin{center}
\includegraphics[width=.9\linewidth]{./pic/backups_plans_20210422_102924.png}
\end{center}

\begin{center}
\includegraphics[width=.9\linewidth]{./pic/backups_plans_20210422_103112.png}
\end{center}

\begin{center}
\includegraphics[width=.9\linewidth]{./pic/backups_plans_20210422_103008.png}
\end{center}

小时候，爸妈因为一次爸爸的不当行为而常年吵架，年幼的我无所适从，本能地认定爸爸是错的爸爸错了！情感上只同妈妈亲。

\begin{center}
\includegraphics[width=.9\linewidth]{./pic/backups_plans_20210422_114555.png}
\end{center}

而且经常哭，我的耳朵总是病着，像是得了顽疾，偶尔好一点儿，但几年时间里就总不断根、总好不了！

长大后的我想，我后来总是胆小怕事、动辙被吓得魂飞魄散、心惊肉跳，可能与那段时期的成长相关联吧。

\begin{center}
\includegraphics[width=.9\linewidth]{./pic/backups_plans_20210422_101017.png}
\end{center}

小学时候，我也说过一次谎话、一堆谎话。

\begin{center}
\includegraphics[width=.9\linewidth]{./pic/backups_plans_20210422_101254.png}
\end{center}

等到大姐到了相亲、谈恋爱的年龄，我就感觉爸妈更喜欢大姐！

\begin{center}
\includegraphics[width=.9\linewidth]{./pic/backups_plans_20210422_074914.png}
\end{center}

而回到当时小学的课堂上，对于自己喜欢的青梅竹马小伙伴，那个时候的我总是会想要捉弄他一下！

\begin{center}
\includegraphics[width=.9\linewidth]{./pic/backups_plans_20210422_100606.png}
\end{center}

小学时候我还是很贪玩儿的，因为我还太小，爸妈也不在学习上给我任何压力，我的成绩总是轻轻松松考个班上前三名就可以了。

\begin{center}
\includegraphics[width=.9\linewidth]{./pic/backups_plans_20210422_100644.png}
\end{center}

小学时候的体育好吗？不喜欢动。

\begin{center}
\includegraphics[width=.9\linewidth]{./pic/backups_plans_20210422_100732.png}
\end{center}

小学官衔：小组长，常年小组长，从一年级到六年级每年都是小组长。

\begin{center}
\includegraphics[width=.9\linewidth]{./pic/backups_plans_20210422_080215.png}
\end{center}

多少年以后的后来，当我真正喜欢上自己那亲爱的表哥，我也曾经潜意识里想要捉弄过表哥！

\begin{center}
\includegraphics[width=.9\linewidth]{./pic/backups_plans_20210422_100000.png}
\end{center}

\begin{center}
\includegraphics[width=.9\linewidth]{./pic/backups_plans_20210422_074743.png}
\end{center}

我想过很多次，为什么小学时候，我没有积极主动地要求在家里晚上要看电视剧，第二天上下放学的路上可以同小伙伴们一起讨论电视剧？

为什么从小学起、感觉整个小学生涯我的话就很少，像是个不怎么说话的哑巴娃娃？

对青梅竹马小男孩的喜欢，于我是一种沉浸，是一种浸泡在里面六年的小女孩情愫。

\begin{center}
\includegraphics[width=.9\linewidth]{./pic/backups_plans_20210422_112117.png}
\end{center}

\begin{center}
\includegraphics[width=.9\linewidth]{./pic/backups_plans_20210422_112135.png}
\end{center}

\begin{center}
\includegraphics[width=.9\linewidth]{./pic/backups_plans_20210422_112152.png}
\end{center}

与自己小学六年青梅竹马兼之同行的，是亲情里、爸妈情感结构里，因外公外婆重男轻女而对我老家舅舅家我的亲表哥格外宠爱溺爱的亲情投射，我同样也是在暗恋着自己的表哥！

如果说青梅竹马只是自己小学六年里发生过、发生着的事，那我对自己亲表哥的暗恋，又持续了多久、持续到什么时候了呢？

这还真是一个放养的、沉浸浸泡式成长呢？！

如果说那段成长是周遭亲人无心之失、无意中形成；

那么多年以后，当我来到异国他乡，当舅舅和我亲爱的表哥都分别播打911之后，我再一次地、这一次被放养到了激烈的竞争环境中。

\begin{center}
\includegraphics[width=.9\linewidth]{./pic/backups_plans_20210422_121139.png}
\end{center}

那么，我接下来的命运，是会在沉默中暴发，抑或是就此沉沦？这是后话。

\begin{center}
\includegraphics[width=.9\linewidth]{./pic/backups_plans_20210422_120028.png}
\end{center}

后来，成长过程中，生活无意中被我们撞见、无形中观察到的事实总是会一再敲响心门：提醒我们自己，我们到底是什么样的人、我们想要成为什么样的人，又或者我们想要寻找怎样的幸福。 

\begin{center}
\includegraphics[width=.9\linewidth]{./pic/backups_plans_20210422_114817.png}
\end{center}

当后来的某天，当我曾经的小伙伴（自己的传记里把他称作了“小伙”）某次偶然间，把他的头、把他的耳朵直接递给我的时候，我当时不免心里又惊了一下。


\section{成长的故事 -- 我和表哥}
\label{sec:org7a2747b}
\begin{itemize}
\item 2011年11月4日，当三大中文媒体对我的人肉已经伤及我自身生活，我必须站出来澄清自己, in Part 1, （San Jose, CA）；

\begin{center}
\includegraphics[width=.9\linewidth]{./pic/dreamer1.png}
\end{center}
\item 4/19/2012 - 6/17/2012, in Part 1, 第二次写至统计专业OPT实习结束（San Jose, CA）；

\begin{center}
\includegraphics[width=.9\linewidth]{./pic/dreamer2.png}
\end{center}
\item 2014年夏天，写于SJSU Library (San Jose State University Public Library, San Jose, CA)

\begin{center}
\includegraphics[width=.9\linewidth]{./pic/dreamer30.png}
\end{center}
\item 2/13/2015 - 12/17/2015(?, Moscow, ID; either and or not San Jose State University Public Library, San Jose, CA)

\begin{center}
\includegraphics[width=.9\linewidth]{./pic/dreamer3.png}
\end{center}

\item I will reorganize the four pdfs, and emphasize keys issues and situations of the whole process, while at the same time to help major population understand what's going on, and what's inside opinions. 虽然这个成长的故事系列是以2011年当三大中文网站（mitbbs.com, wenxuecity.com and backchina.com）中文媒体对我的人肉与网上评论伤及我的正常生活时，我站出来开始写自己的自传，并分四次在四个不同的时间段，不同舆论或事件压力下或是网上澄清，或是网上求助以便能帮我泄掉一部分当时自己的压力，分四次于不同的地点纪录了的自己的主要生活，纪录到2015年计算机硕士学位结束。
\item 这一次，这里，我会以事件主要人物及其相关主要事迹的人物列传、或/和大事记、大冲突记的形式来重新组织语言，重述我的整个成长史与大事记、大冲突记，来帮助自己成长、并帮助社会大众认清事情所有环节真相的目的。但鉴于时间有限，我会以剧情梗概的形式每天大致纪录与一个相关人物某件或某几件事的进展、或一天一两个主要事件，并将已经完成了的四个部分作为原始事件纪录的细节参考供索引，并争取做到每日更新一篇，到我把先前与这个教授舅舅的所有冲突的这件事情具体讲述清楚，以供大家共同去探讨事情的真相到底如何，有一个更能为大家所接受或理解的底层社会小人物的心灵成长史。
\end{itemize}

\section{重返校园}
\label{sec:orgd2fbadd}

\begin{center}
\includegraphics[width=.9\linewidth]{./pic/backups_plans_20210414_161755.png}
\end{center}

\begin{center}
\includegraphics[width=.9\linewidth]{./pic/backups_plans_20210414_161857.png}
\end{center}

\begin{center}
\includegraphics[width=.9\linewidth]{./pic/backups_plans_20210414_161940.png}
\end{center}

如同2014年夏天那第三次地站出来写自己的传记般，2012年的夏天，在5月底结束了那份统计OPT的最后的三个月的工作后，我重新返校了，去从头开始读一个计算机专业的硕士。

\begin{center}
\includegraphics[width=.9\linewidth]{./pic/backups_plans_20210419_103028.png}
\end{center}

具体的我是什么时候与学校取得联系，并快速地申请了计算机专业，我已经想不起来，无法追忆了。我应该是6月份、7月份还住在加州的（7月底8月头回得学校？），根据系里小秘建议和提供的联系方式，我 \textbf{当天} （我昨天读到这个字，把自己读哭了！）就与当时系里帮我分配的导师取得了联系，并就秋季选课的事情与导师协商、讨论。

为什么当时的自己就那么迫切地想要与系里为我分配的导师、甚至于还没有见过面的导师，去讨论还远在一两个月之后的自己读计算机专业的选课问题呢？

因为我不够独立，我有依赖性，我还不够自信。

\begin{center}
\includegraphics[width=.9\linewidth]{./pic/backups_plans_20210419_103828.png}
\end{center}

你看，在先前的要不要读一个计算机专业的时候，我第一时间写邮件征询我亲爱的表哥与舅舅的意见，我的表哥没有理我，舅舅也只给了我四个字“We have no suggestions.”

\begin{center}
\includegraphics[width=.9\linewidth]{./pic/backups_plans_20210419_104129.png}
\end{center}

在一年前的7月份，因为朋友的怂勇我写邮件向表哥表达过结婚意愿后，舅舅在邮件里警告我，舅舅在邮件里对我使用冷暴力！我的自尊心受到了极大的伤害，一旦我有了工作、有了维持维护自己尊严的工作（8月头），我便正式工作开始之前就怒气冲冲地杀回去找舅舅报仇了，还惹得舅舅真的播打了911！

\begin{center}
\includegraphics[width=.9\linewidth]{./pic/p1p34.png}
\end{center}

\begin{center}
\includegraphics[width=.9\linewidth]{./pic/backups_plans_20210419_104535.png}
\end{center}

如果说2008年寒假从加州回到学校的我给舅舅写邮件，表达了我那次去加州，因为时间紧急，没有机会没能帮舅舅带任何礼物回来的疚意，舅舅回复我的邮件曾经说过的两个字“Welcome home.”曾经深深地感动过那些年月里的我！

\begin{center}
\includegraphics[width=.9\linewidth]{./pic/backups_plans_20210419_105423.png}
\end{center}

那么这次舅舅用更长的邮件、两倍的字数——四个字对我征求意见的回复，让那个受到过舅舅的冷暴力警告、并在接下来的一两个星期内杀回家去找舅舅报过仇、并且舅舅真的播打了911的自己，真正感觉到了我最亲爱的表哥、这我在美国再一次地找上门去相认才得到的我的阔别10年的舅舅（第一次认舅舅是在国内，1997年暑假的时候），虽然表哥和舅舅都是我的远亲、但他们在我这里、在我的世界里却是血浓于水、至关重要、永远也不想割舍的亲情，正在慢慢离我远去、渐行渐远！

在接下来远近一年、大半年的时间里，我反复体会着、咀嚼着那份亲情远离的深深痛楚！

\begin{center}
\includegraphics[width=.9\linewidth]{./pic/backups_plans_20210419_113045.png}
\end{center}

\begin{center}
\includegraphics[width=.9\linewidth]{./pic/backups_plans_20210419_113136.png}
\end{center}

\begin{center}
\includegraphics[width=.9\linewidth]{./pic/backups_plans_20210419_113202.png}
\end{center}

舅舅警告和真正亲自播打了911的当时——那时那会儿，我就不会痛吗？痛——是一定的！在当时，痛的表现形式是彻底割舍：我想我只要做好自己、努力工作，忘掉表哥，我就能走进自己的新时代！

但这份痛的深远影响却留在了接下来的反刍、迷失与找回自己的岁月里。 

\section{重返校园（2）}
\label{sec:orgea8824c}

\begin{center}
\includegraphics[width=.9\linewidth]{./pic/backups_plans_20210420_115754.png}
\end{center}

去年、今年的统计29个月OPT期间，舅舅和表哥先后播打了911期间，我以为舅舅的冷暴力播打911后，我以为我是不痛的，因为我转身就要走向自己的新时代了！11年8月当舅舅真正播打了911之后，我想，我只要做好自己、努力工作、忘掉表哥，我就能走进自己的新时代！

\begin{center}
\includegraphics[width=.9\linewidth]{./pic/backups_plans_20210420_120854.png}
\end{center}

当年的自己，2009年秋季学期，因为对系里一位漂亮、打份相对前卫的美女老师的不信任，我压根儿就不敢跟她作研究！现在，系里为我分配的这个导师，我就熟吗？我就敢吗？可为什么她就是那么迫切地想要与他联系呢？

直到我这次重新回读、回味和对比、对照着自己这些年的成长来写回忆录，被当年邮件里的那一个字读哭，禁不住叹喟当年的那个孩子！

2012年的事情，过去快9年了，好多事情、故事以及细节都被自已遗忘了。所以这两天再回去读（今年三月之前、至少15年之后，从来不曾回去重新读起过！），还是会常常把自己读哭的。

\begin{center}
\includegraphics[width=.9\linewidth]{./pic/backups_plans_20210420_114702.png}
\end{center}

在我向导师介绍了自己，表达需要选课诉求后，导师首先问我，你的目标是什么？

\begin{center}
\includegraphics[width=.9\linewidth]{./pic/backups_plans_20210419_084838.png}
\end{center}

但当时的我，对于导师提出来的这个问题，我是没有明确目标或者说专业领域的方向的，因为我不熟不懂！

如果说心里有相对明确的人生目标，我想还是应该是比较喜欢实习期间的那些工作环境（希望将来能够工作），每天能够激情飞扬地完成一天的工作，晚上下班后便再没有了工作上的压力与顾虑，每天晚上回到家都可以安安稳稳地睡个好觉 。可是，这，好像不是导师想问的问题。

他问的应该是研究的兴趣、科研的方向？可是为什么我会想要走科研的道路呢？这应该是当时的情商弱弱读不出来的潜在问题了。 

\begin{center}
\includegraphics[width=.9\linewidth]{./pic/backups_plans_20210420_121822.png}
\end{center}

导师问及我的编程经验，我便回忆、向导师一一列举了我所有的编程相关的课程与经验。

\begin{center}
\includegraphics[width=.9\linewidth]{./pic/backups_plans_20210419_085025.png}
\end{center}

以前的成绩单：

\begin{center}
\includegraphics[width=.9\linewidth]{./pic/backups_plans_20210419_095006.png}
\end{center}

\begin{center}
\includegraphics[width=.9\linewidth]{./pic/backups_plans_20210419_093849.png}
\end{center}

\begin{center}
\includegraphics[width=.9\linewidth]{./pic/backups_plans_20210419_093428.png}
\end{center}

\begin{center}
\includegraphics[width=.9\linewidth]{./pic/backups_plans_20210419_093456.png}
\end{center}

《计算机程序语言设计》：3个学分。《计算机基础》的1个学分因为我补考才过的，没有学分。

\begin{center}
\includegraphics[width=.9\linewidth]{./pic/backups_plans_20210420_122207.png}
\end{center}

说我对这个专业带着“敬畏”，也是因为当年99年春天的第二学期计算机基础课上机考试，我有一个什么地方没有弄好，程序没能保存下来，结果那门课我被要求补考过（学分还记成了是0个学分，原本我应该是拿到1个学分）。那是整个上学期间（学生生涯？）唯一一次补考。（叹一下：放养、同时又以小混混为楷模长大的孩子、一切的重大成长，都以痛苦深刻的教训当拌脚石来推动促进成长，成长得好痛苦、好悲催！）

这里也顺带提一句：我的《成长的故事》写到此，绝大部分的读者都已然清楚，我原本高考没有考好，所以上大学选择了当初舅舅帮忙建议我上我的农林院校。来到美国后，在语言有困难的情况下，舅舅帮忙经济担保我读《统计》的硕士，而现在我想要顺应自己的兴趣去探索的是《计算机》，想拿计算机的硕士学位。这在国内教育体制下是非常困难的。

因为高考考完之后，我没能去想、也可能上了大学后也是没有足够的勇气去放弃、并重回高三去复读，以期待重新考取更感兴趣或更有前途的专业，那么在国内当时的教育体制下，我人生最大的不幸——高考没考好所导致的这个农林院校的专业就很有可能、将会跟随我一辈子，如影随形。

高考之后，农家孩子学业的道路上，我们可以再重新选择专业的机会就只有研究生入学考试，但如果选择转专业，并且是通过研究生入学考试这样一项硬指标来作为唯一评判标准，对于非专业、非科班出生的考生或门外汉（比如我农林院校的本科书，想要考研究生并想同时转成读计算机专业硕士）来说，从获胜希望上、竞争激烈程度上来说，都是一种致命的打击。因为我们我们作为人的本能的个人兴趣，在强大的以考试成绩为唯一标准、与受过四五年大学本科科班教育的本专业考生相比，在强大的国家选拔机制国家机器的运转面前，我们个人的那一点儿兴趣、因为爱好喜欢而迈出的微尘一小步，是多么地渺小、微不足道、不值一提，在强硬的选拔机制面前，那微尘一小步，压根儿就不会再有任何的舞动空间！

所以，我们就成为了模式化教育长大的克隆人。而最终成就不同克隆人之间区别的就成为了：他们的成长环境与所成就的个性、他们学习工作的竞争力与学习工作环境的相系制约，一如我——《成长的故事——我和表哥》的自传作者，现在所想要讲述的，除了我这亲爱的表哥与我——这终将浸透岁月的爱情，同时讲述的，也就包括了我——一个克隆人的心灵成长史与国家考试选拔机制、学习工作环境与竞争机制的相互制约、相互作用等。

这个克隆人没有望穿、透视浩瀚星空的透彻与洞察力，仅以微尘之眼观察环绕着她的这个周围的世界。

\section{重返校园（3）}
\label{sec:orgace0025}

（一） 学习目的

\begin{center}
\includegraphics[width=.9\linewidth]{./pic/backups_plans_20210421_123440.png}
\end{center}

在系里小秘给了我系里为我安排的导师的“当天”，在写给自己导师的第一封邮件里，我向自己的导师列具了自己目前计划和钟意想要选的课程。

\begin{center}
\includegraphics[width=.9\linewidth]{./pic/backups_plans_20210421_123603.png}
\end{center}

在写给自己导师的第一封邮件里，当初的自己，也表达了对于来年暑假可能会有的实习，自己心里的想法与想要征求的意见建议。

这转专业转向、步入计算机专业之初的想往，后来实现了吗？这是后话。 

\begin{center}
\includegraphics[width=.9\linewidth]{./pic/backups_plans_20210421_125334.png}
\end{center}

在我向导师介绍了自己，表达需要选课诉求后，导师首先问我，你的目标是什么？这个学习的目标，应试是狠重要的！

\begin{center}
\includegraphics[width=.9\linewidth]{./pic/backups_plans_20210419_084838.png}
\end{center}

但当时的我，对于导师提出来的这个问题，我是没有明确目标或者说具体的专业领域里的方向的，因为我不熟不懂！

如果说心里有相对明确的学习这个计算机专业之后的人生目标，我想还是应该是比较喜欢实习期间的那些工作环境（希望将来能够工作），每天能够激情飞扬地完成一天的工作，晚上下班后便再没有了工作上的压力与顾虑，每天晚上回到家都可以安安稳稳地睡个好觉 。可是，这，好像不是导师想问的问题。

他问的应该是研究的兴趣、科研的方向？可是为什么我会想要走科研的道路呢？这应该是当时的情商弱弱读不出来的潜在问题了。 

当时的我没有读懂导师的问题，也没有深入地去思考导师想要问的究竟是什么，以至于这个最重要的问题，仿佛被自己华丽丽地忽视了？！！！

而这个问题狠重要！当时邮件中的自己忽略了，回到学校里去的我，见到自己的导师后，应该再就这个问题具体与导师再好好讨论一次，好像才对。

（二） 编程基础

\begin{center}
\includegraphics[width=.9\linewidth]{./pic/backups_plans_20210421_130001.png}
\end{center}

导师也问过我的编程经验，因为前一篇思路不够清晰、略有涉及、这篇略过不再冗述。

\begin{center}
\includegraphics[width=.9\linewidth]{./pic/backups_plans_20210421_125854.png}
\end{center}

导师所阐述的他对于这个计算机专业、教书育人的基本理念。

\begin{center}
\includegraphics[width=.9\linewidth]{./pic/backups_plans_20210421_125909.png}
\end{center}

以及他再具体一点儿、涉及到CS121这门课，我选与不选的区别与影响。

\begin{center}
\includegraphics[width=.9\linewidth]{./pic/backups_plans_20210419_085149.png}
\end{center}

导师的这一理念我也是同意的，一开始我也就肯定过。

（三） 专业兴趣与方向

\begin{center}
\includegraphics[width=.9\linewidth]{./pic/backups_plans_20210421_090244.png}
\end{center}

\begin{center}
\includegraphics[width=.9\linewidth]{./pic/backups_plans_20210421_090620.png}
\end{center}

当导师为我指明生物信息学这个方向，比较适合农林院校专业、有着统计背景、又很想要学习一点儿一些计算机的我时，我是欣然接受、还是心蒙尘埃、本能规避呢？

\begin{center}
\includegraphics[width=.9\linewidth]{./pic/backups_plans_20210421_090714.png}
\end{center}

\begin{center}
\includegraphics[width=.9\linewidth]{./pic/backups_plans_20210419_094833.png}
\end{center}

\begin{center}
\includegraphics[width=.9\linewidth]{./pic/backups_plans_20210419_094910.png}
\end{center}

\begin{center}
\includegraphics[width=.9\linewidth]{./pic/backups_plans_20210419_085406.png}
\end{center}

我对导师讲述了我在国内的硕士学位时在中国农业科学院蔬菜花卉研究所已经略微接触和学习过一门《生物信息学》的课程，但因为我只有极其有限的计算机基础，我并不清楚导师所指出的这个方向是否就将成为我最感兴趣、最想要学习的方向、走的路。但因为这扇窗户是为自己敞开着的，我这个学期可以试着选一门这样的课先听听看。

\begin{center}
\includegraphics[width=.9\linewidth]{./pic/backups_plans_20210421_090244.png}
\end{center}

\begin{center}
\includegraphics[width=.9\linewidth]{./pic/backups_plans_20210419_085322.png}
\end{center}

如果说之前我还有想要留下来，征服一门自己曾经非常感兴趣的专业，那么在导师如此肯定绝决的语气下，我已是心生退意，因为我没有那么多的钱来读完这样一门专业。

以前与导师的通信基础上收到他的邮件，最迟第二天就回了，我次，我等了四天才回复他的邮件。 

\begin{center}
\includegraphics[width=.9\linewidth]{./pic/backups_plans_20210421_090714.png}
\end{center}

\begin{center}
\includegraphics[width=.9\linewidth]{./pic/backups_plans_20210421_091059.png}
\end{center}

第二遍肯定：与导师可以达成共识的地方：计算机的编程基础非常重要——非常重要！

如果我对自己的编程没有足够的信心、如果我不敢跳过CS121，那么我应该先在第一学期选择这样一门课，而把其它的课程往后排靠后。

\begin{center}
\includegraphics[width=.9\linewidth]{./pic/backups_plans_20210421_091227.png}
\end{center}

因为如果我这学期选最基础的编程课CS121，那么接下来的一个学期我是同样可以选择200层次的课程的，而不是要等待一年到来年的秋季学期。

\begin{center}
\includegraphics[width=.9\linewidth]{./pic/backups_plans_20210421_091430.png}
\end{center}

我清楚地表达了我想要这个学期多选一点儿课的原因：我的学费不够，所以我希望我们能够有一个更好的选课方案（让我能够在有限的学习经费下顺利地完成我的学业）。

\begin{center}
\includegraphics[width=.9\linewidth]{./pic/backups_plans_20210419_084632.png}
\end{center}

我舍不得表哥，与留下读一个计算机专业可能会有的学费不够的问题下，朋友的表态还是给了我底气：我留下来，一定会有解决办法的！

\begin{center}
\includegraphics[width=.9\linewidth]{./pic/backups_plans_20210421_091700.png}
\end{center}

当前一封邮件里老师对我的计算机基础知识、编程基础存在怀疑的时候，我是这样表达自己对于编程的兴趣或自己的长处与爱好的：

\begin{center}
\includegraphics[width=.9\linewidth]{./pic/backups_plans_20210419_085536.png}
\end{center}

那么当我表达了我对编程的强烈兴趣——它是我对这个计算机专业最大的兴趣，列出了我的所有计算机相关的基础，导师的反应又会是怎样的呢？

\begin{center}
\includegraphics[width=.9\linewidth]{./pic/backups_plans_20210421_092026.png}
\end{center}

\begin{center}
\includegraphics[width=.9\linewidth]{./pic/backups_plans_20210421_092635.png}
\end{center}

当我前一封邮件等了四天才回复他，导师也把对我的回复邮件放到了四天之后。 

基础计算机技能的重要性。关于学位，当时的自己应该又是把这句话给华丽丽地读丢了，所以后来有了好玩儿的导师以为我不要学位的后续故事事件。 

\begin{center}
\includegraphics[width=.9\linewidth]{./pic/backups_plans_20210421_092842.png}
\end{center}

导师再次强调编程的重要性，因为我读的是硕士学位，B的成绩对于导师来说是个问题、惊叹号、是不可以接受的！

导师的态度似乎是强硬的，但是导师过分了吗？

我们回过头来看看我写给导师的第一封邮件、自己的介绍信里，所曾经表达过的对于自己想要选择课程的某些想法：

\begin{center}
\includegraphics[width=.9\linewidth]{./pic/backups_plans_20210420_114525.png}
\end{center}

读到这里，那么我也想要问：你到底是为什么自己曾经上过的课程，你还想要再学一遍，而不是去选择其它课程、学习新的知识？

这个编程的基础，你到底是有，还是没有？你对自己的编程，到底是有信心，还是没有足够的信心？

转到这个计算机专业、想要去读计算机的信心，你到底是有，还是没有呢？

\begin{center}
\includegraphics[width=.9\linewidth]{./pic/backups_plans_20210421_134744.png}
\end{center}

\begin{center}
\includegraphics[width=.9\linewidth]{./pic/backups_plans_20210421_134614.png}
\end{center}

99年第二学期学《计算机基础》，因为上机考试最后的文件我没能保存好，我补考过，没拿到学分，所以从第三学期才好好学习这门课，成绩也是真的很好！

\begin{center}
\includegraphics[width=.9\linewidth]{./pic/backups_plans_20210421_135356.png}
\end{center}

\begin{center}
\includegraphics[width=.9\linewidth]{./pic/backups_plans_20210421_135421.png}
\end{center}

《数据库原理与应用》课上得也还可以，比《生物信息学》上得好。主要是那会儿生物信息学的课安排在每周哪一两天中午下午两三点钟上，感觉那个学期头好昏啊，好多时间都不知道老师在讲什么！

\begin{center}
\includegraphics[width=.9\linewidth]{./pic/backups_plans_20210421_093121.png}
\end{center}

这个对计算机专业常怀敬畏之心，有一定的编程基础、但成绩得的是B的学生，这个时候，被导师几乎绝决的语气吓到，犹豫起来！

\begin{center}
\includegraphics[width=.9\linewidth]{./pic/backups_plans_20210421_093412.png}
\end{center}

\begin{center}
\includegraphics[width=.9\linewidth]{./pic/backups_plans_20210419_085713.png}
\end{center}

这次，再等十天之后才回复给导师的邮件，我清楚地解释了2009年秋天我统计最后一个学期选修一门计算机课CS120只得了个B的原因，向导师清楚地表达着自己对计算机、或者更确切地说，
对编程的兴趣。

\begin{center}
\includegraphics[width=.9\linewidth]{./pic/backups_plans_20210421_134255.png}
\end{center}

我向导师清楚地解决自己迟疑十天、犹豫的原因是，农家穷孩子对自己自费回来读这样一个专业经济上的担心。

当时的我，没能想得长远的是，我既然想要留下来读计算机专业，如果我的真的如对自己认定的这般有着很好的分析能力、与解决问题的能力、喜欢编程，那么等我真正学得好、表现得好，真正能够表现得还比较突出的时候，我还是有机会、可能能够从系里、院里获得经济资助的，一如后来我有拿到一个学年的奖学金，这是后话。 

如果我们对一门专业有着强烈的兴趣与爱好，如果与本科生本专业同科班出身的他们相比，我们那仅只因为兴趣、爱好而迈出的微尘一小步尚不足以与他们抗衡的初始，在我们没有任何竞争力的初始，为了摆脱掉一个如影随形、可能今生都终将跟随自己的尚不够有兴趣的专业，为了自己心底的兴趣与爱好、那么我们自费去学习一个新的自己有兴趣的专业，为了自已的兴趣爱好与今后工作上的满足感与成就感，我们拿自己的钱、自己今生工作的积蓄，去与那个尚且不够成熟、尚且不够定性的自己打一场翻身仗，成功与否，我们不是也在我们还相对年轻的年龄、为了我们自己、轰轰烈烈地活过一场，我们对得起自己、对得起我们的青春年华、没有遗憾，不是吗？

这一次，导师清楚地读出了自己邮件里所表达出的担心，我的导师这次没有像上次一样，没有再与我先前等同样的四天、或是等这次的十天再回复我的邮件，而是“当天”就给我回复了他的态度与立场：

\begin{center}
\includegraphics[width=.9\linewidth]{./pic/backups_plans_20210421_093907.png}
\end{center}

\begin{center}
\includegraphics[width=.9\linewidth]{./pic/backups_plans_20210419_085846.png}
\end{center}

导师也为我解释了，他上一次话说得重、或是他摆明当时态度的原因：医生看病是要根据征状来的，而我也是就你曾经的成绩就事说事！应该也是希望我能够理解他曾经的立场与态度吧！

至此，我们终于是达成共识：我是有选课余地的！我可以选CS121，或是跳过这门基础编程课，全凭、完全根据我自己的信心与信念！

\begin{center}
\includegraphics[width=.9\linewidth]{./pic/backups_plans_20210421_094226.png}
\end{center}

\begin{center}
\includegraphics[width=.9\linewidth]{./pic/backups_plans_20210419_085907.png}
\end{center}

至此，当年的弱弱也已经清楚地认识到，与这个从未谋面的系里为我分派的导师，半个月内简短的几封邮件，已经消除了我选课与将来毕业精神上的压力。我根据自己的基础与兴趣，我有选、与不选某门课、某些课的余地！同样的，我学得好，我也就会能够顺利地毕业，我大可不必过于担心！

而我这个肉肉的、肉奶奶般行动缓慢的女生，并没能如自己的导师理解了我的担心与忧虑、“当天”就回复了我的邮件般，再“当天”把邮件回复回去，而是第二天才回复了导师的邮件，并对导师与自己之间所达成的选课余地表达感谢，虽然当时的自己并没有明说：这个秋季转读《计算机专业》第一学期我到底是选、还是不选这门最基础的C++ CS121 编程课！

\begin{center}
\includegraphics[width=.9\linewidth]{./pic/backups_plans_20210419_085958.png}
\end{center}

系里大牛帮我列出的选课计划，我仍然是没有安全感的，于是回到导师的办公室，导师再重新帮我列一个选课计划！

\begin{center}
\includegraphics[width=.9\linewidth]{./pic/backups_plans_20210421_141850.png}
\end{center}

转专业第一学期选课的结果与过程。

\begin{center}
\includegraphics[width=.9\linewidth]{./pic/backups_plans_20210421_141921.png}
\end{center}

真正试上过一两个周的课后，我自动把自己导师的过于严格的课换掉了、换成了其它课程。

读者，你看，与自己的导师早早地联系，那么多邮件过后，我还是选择了自己认为最对的课程，在自己导师已经允许我跳过CS121编程课的前提下自己坚持仍然选课了，并根据自己的状态换掉了自己导师过严的系统课：一切都仍以自己的需要与状态作为选与不选某门课、换与不换某门课的标准，仍然是走心派呀，那这个导师全程属于打酱油角色？

那么，让我们回到前天把自己读哭的那个“当天”，为什么我就那么迫切地想要、需要与这个系里为我分配的、尚未见过面的陌生的导师取得联系、并讨论远在一两个月之后的秋季学期的选课呢？

能够如当年1997年认舅舅般本能地冲上去、如同10年12月与表哥相处几天的我找表哥告别时会主动向表哥求抱抱、索求拥抱一样，当时的自己——如同那天地之间一个孤独的孩子，能够与系里自己将来的导师及时地建立起联接，她的世界就会多一分的力量与坚强！

回到学校来读书，我就又来到了表哥旁边的学校，又重新回到表哥的身边。我亲爱的表哥，他这段时间会是在做些什么呢？


\begin{center}
\includegraphics[width=.9\linewidth]{./pic/backups_plans_20210421_142818.png}
\end{center}

表哥手里拿着的粉红色口袋！我当时的心理想的应该是，表哥一个40岁的大男人，拿着个粉色口袋是要送给哪家姑娘？

2020年夏天，当我去做instacart帮别人购物，有一次一个客户要我从家旁边的H-mart买洗菜的盆子，并注明提醒：No pink please!我心想，是什么人不喜欢粉红色呢？界时，从2018年我回国探亲开始，我在家里已经攒下了各种各样的粉红色保温杯、玻璃水杯和粉红色袖珍电饭煲等。我都不明白为什么我攒下了这么多粉红色的器具？

到现在写回学校读书这几篇文的这两天，我才、才、才、才想明白！！！

\begin{center}
\includegraphics[width=.9\linewidth]{./pic/backups_plans_20210421_095520.png}
\end{center}

亲爱的表哥，你是从一开始就知道：我身披铠钾、全副武装、一副战威威、军功威赫、沙场战将的样子，内心里却是始终住着个少女心小弱弱的么？麽麽嗒！！！
\end{document}